\documentclass[11pt, a4paper]{article}

\usepackage[margin=2.5cm]{geometry}
\usepackage{listings}
\usepackage{xcolor}
\usepackage{hyperref}
\usepackage{titlesec}
\usepackage{enumitem}
\usepackage{parskip}
\usepackage{microtype}
\usepackage{booktabs}
\usepackage{fancyhdr}

% ── Color palette ──────────────────────────────────────────────────────────────
\definecolor{codebg}{HTML}{F5F5F5}
\definecolor{codeframe}{HTML}{CCCCCC}
\definecolor{kwcolor}{HTML}{0000CC}
\definecolor{strcolor}{HTML}{CC0000}
\definecolor{cmtcolor}{HTML}{007700}
\definecolor{numcolor}{HTML}{AA00AA}
\definecolor{builtincolor}{HTML}{CC6600}
\definecolor{decocolor}{HTML}{AA007F}
\definecolor{accentblue}{HTML}{1A1A6E}

% ── listings: Python style ─────────────────────────────────────────────────────
\lstdefinestyle{python}{
  language=Python,
  backgroundcolor=\color{codebg},
  frame=single,
  rulecolor=\color{codeframe},
  framesep=6pt,
  xleftmargin=10pt,
  xrightmargin=10pt,
  basicstyle=\ttfamily\small,
  keywordstyle=\color{kwcolor}\bfseries,
  stringstyle=\color{strcolor},
  commentstyle=\color{cmtcolor}\itshape,
  numberstyle=\tiny\color{gray},
  numbers=left,
  numbersep=8pt,
  stepnumber=1,
  showstringspaces=false,
  breaklines=true,
  breakatwhitespace=true,
  tabsize=4,
  captionpos=b,
  morekeywords={async, await, True, False, None},
  emph={@mcp.tool, @mcp.resource, @mcp.prompt},
  emphstyle=\color{decocolor}\bfseries,
  literate={->}{{\texttt{->}}}{2},
}

\lstset{style=python}

% ── hyperref ───────────────────────────────────────────────────────────────────
\hypersetup{
  colorlinks=true,
  linkcolor=accentblue,
  urlcolor=accentblue,
  citecolor=accentblue,
  pdftitle={mcp-modular Developer Guide},
  pdfauthor={mcp-modular},
}

% ── headers / footers ──────────────────────────────────────────────────────────
\pagestyle{fancy}
\fancyhf{}
\fancyhead[L]{\small\textit{mcp-modular}}
\fancyhead[R]{\small\textit{Developer Guide}}
\fancyfoot[C]{\thepage}
\renewcommand{\headrulewidth}{0.4pt}

% ── title ──────────────────────────────────────────────────────────────────────
\title{%
  {\Large\bfseries mcp-modular}\\[4pt]
  {\large Developer Guide}
}
\author{}
\date{}

% ───────────────────────────────────────────────────────────────────────────────
\begin{document}
\maketitle
\thispagestyle{fancy}

\tableofcontents
\newpage

% ══════════════════════════════════════════════════════════════════════════════
\section{Overview}
% ══════════════════════════════════════════════════════════════════════════════

\textbf{mcp-modular} is a modular \href{https://modelcontextprotocol.io}{Model Context Protocol (MCP)} server built with \href{https://github.com/jlowin/fastmcp}{FastMCP}. It pairs with a terminal chat client backed by \href{https://ollama.com}{Ollama} to provide a local, extensible AI assistant that can call tools, serve resources, and inject structured prompts.

The defining characteristic of this project is \textbf{automatic module discovery}: you drop a new Python file into the right folder and the server picks it up on the next start — no registration, no imports to edit.

% ══════════════════════════════════════════════════════════════════════════════
\section{Prerequisites}
% ══════════════════════════════════════════════════════════════════════════════

\begin{itemize}[leftmargin=2em]
  \item Python 3.11 or later
  \item \href{https://ollama.com/download}{Ollama} installed and running locally
  \item The \texttt{llama3.1:8b} model pulled in Ollama: \texttt{ollama pull llama3.1:8b}
\end{itemize}

% ══════════════════════════════════════════════════════════════════════════════
\section{Setup}
% ══════════════════════════════════════════════════════════════════════════════

\subsection{Clone and install}

Create a virtual environment, activate it, and install the dependencies declared in \texttt{requirements.txt}:

\begin{lstlisting}[language=bash, style=python, numbers=none]
git clone <repo-url> mcp-modular
cd mcp-modular

python -m venv venv
source venv/bin/activate          # Windows: venv\Scripts\activate

pip install -r requirements.txt
\end{lstlisting}

The only runtime dependencies are:

\begin{center}
\begin{tabular}{ll}
\toprule
Package & Version \\
\midrule
\texttt{mcp} & 1.26.0 \\
\texttt{ollama} & 0.6.1 \\
\bottomrule
\end{tabular}
\end{center}

\subsection{Start the server}

\begin{lstlisting}[language=bash, style=python, numbers=none]
python server.py
\end{lstlisting}

By default the server listens on \texttt{0.0.0.0:8000} using the SSE transport. You can override this with environment variables:

\begin{center}
\begin{tabular}{lll}
\toprule
Variable & Default & Effect \\
\midrule
\texttt{MCP\_HOST} & \texttt{0.0.0.0} & Bind address \\
\texttt{MCP\_PORT} & \texttt{8000} & TCP port \\
\texttt{MCP\_TRANSPORT} & \texttt{sse} & Transport (\texttt{sse} or \texttt{stdio}) \\
\bottomrule
\end{tabular}
\end{center}

\subsection{Start the client}

In a separate terminal (with the virtual environment activated):

\begin{lstlisting}[language=bash, style=python, numbers=none]
python client.py
\end{lstlisting}

The client connects to \texttt{http://localhost:8000/sse} by default. Adjust with \texttt{MCP\_HOST} and \texttt{MCP\_PORT} as needed.

% ══════════════════════════════════════════════════════════════════════════════
\section{Project Structure}
% ══════════════════════════════════════════════════════════════════════════════

\begin{lstlisting}[language=bash, style=python, numbers=none]
mcp-modular/
|-- mcp_instance.py        # Shared FastMCP instance
|-- server.py              # Server entry point + auto-discovery
|-- client.py              # Interactive chat client
|-- client_helper.py       # Parsing and formatting utilities
|-- requirements.txt
|
|-- tools/                 # @mcp.tool() modules
|   |-- __init__.py
|   |-- coke_product_transparency.py
|   `-- coke_crisis_protocol.py
|
|-- resources/             # @mcp.resource() modules
|   |-- __init__.py
|   |-- coke_advertisements.py
|   |-- coke_campus_locations.py
|   |-- coke_regulations.py
|   `-- data/              # Static data files consumed by resources
|       |-- massachusetts_regulations.json
|       `-- mit_campus.txt
|
|-- prompts/               # @mcp.prompt() modules
|   |-- __init__.py
|   `-- coke_evaluate_beverage.py
|
`-- documentation/
    `-- guide.tex
\end{lstlisting}

\subsection{Folder descriptions}

\begin{description}[leftmargin=3em, style=nextline]
  \item[\texttt{tools/}]
    Each file defines one or more \emph{tools} — callable functions that the LLM can invoke to perform actions or look up data. Decorated with \lstinline{@mcp.tool()}.

  \item[\texttt{resources/}]
    Each file exposes one or more \emph{resources} — read-only content endpoints addressable by a URI. Resources can be static (fixed URI) or templated (URI with \texttt{\{param\}} placeholders). Decorated with \lstinline{@mcp.resource("scheme://path")}.

  \item[\texttt{prompts/}]
    Each file defines one or more \emph{prompts} — reusable message templates that the client can inject into the conversation. Decorated with \lstinline{@mcp.prompt()}.

  \item[\texttt{resources/data/}]
    Static data files (JSON, plain text) loaded at runtime by the resource modules. Not auto-discovered; loaded explicitly by their owning module.

  \item[\texttt{mcp\_instance.py}]
    The single source of truth for the \lstinline{FastMCP} object. Importing from here ensures every module shares the same server instance.

  \item[\texttt{server.py}]
    Discovers and loads all modules, then starts the transport.

  \item[\texttt{client.py}]
    Interactive terminal chat backed by Ollama.

  \item[\texttt{client\_helper.py}]
    Parsing and ANSI-color utilities used by the client (see Appendix \ref{appendix:client_helper}).
\end{description}

% ══════════════════════════════════════════════════════════════════════════════
\section{The Server (\texttt{server.py})}
% ══════════════════════════════════════════════════════════════════════════════

\begin{lstlisting}[caption={server.py}]
import os
import importlib
import pkgutil

import prompts
import tools
import resources
from mcp_instance import mcp

for pkg in [prompts, tools, resources]:
    for _, name, _ in pkgutil.iter_modules(pkg.__path__, pkg.__name__ + "."):
        importlib.import_module(name)


def main():
    transport = os.getenv("MCP_TRANSPORT", "sse")

    if transport == "sse":
        mcp.settings.host = os.getenv("MCP_HOST", "0.0.0.0")
        mcp.settings.port = int(os.getenv("MCP_PORT", "8000"))
        mcp.run(transport="sse")
    else:
        mcp.run(transport="stdio")


if __name__ == "__main__":
    main()
\end{lstlisting}

\paragraph{Auto-discovery loop (lines 7--9).}
\texttt{pkgutil.iter\_modules} walks each package directory and yields the dotted module name of every \texttt{.py} file it finds. \texttt{importlib.import\_module} then imports it. Because every tool/resource/prompt module calls a decorator on \lstinline{mcp} at import time, this single loop is all that is needed to register every capability.

\paragraph{Transport selection.}
Two transports are supported:
\begin{itemize}
  \item \textbf{SSE} (default) — the server opens an HTTP endpoint. Suitable for any client that can speak HTTP; allows the client and server to run in separate processes.
  \item \textbf{stdio} — the server communicates over standard input/output. Useful when embedded inside another process (e.g.\ a code editor plugin).
\end{itemize}

\paragraph{Shared instance (\texttt{mcp\_instance.py}).}
\begin{lstlisting}[caption={mcp\_instance.py}]
from mcp.server.fastmcp import FastMCP

mcp = FastMCP("my-mcp-server")
\end{lstlisting}

Every module imports \lstinline{mcp} from here. This guarantees decorators like \lstinline{@mcp.tool()} all register on the same \lstinline{FastMCP} object the server will start.

% ══════════════════════════════════════════════════════════════════════════════
\section{The Client (\texttt{client.py})}
% ══════════════════════════════════════════════════════════════════════════════

The client is a self-contained terminal application. It connects to the MCP server, queries what is available, and enters an interactive chat loop driven by a local Ollama model.

\begin{lstlisting}[caption={client.py --- startup and main()}]
MODEL = "llama3.1:8b"

async def main():
    host = os.getenv("MCP_HOST", "localhost")
    port = int(os.getenv("MCP_PORT", "8000"))
    url = f"http://{host}:{port}/sse"

    async with sse_client(url) as (read, write):
        async with ClientSession(read, write) as session:
            await session.initialize()

            tools_resp, resources_resp, templates_resp, prompts_resp = \
                await asyncio.gather(
                    session.list_tools(),
                    session.list_resources(),
                    session.list_resource_templates(),
                    session.list_prompts(),
                )

            ollama_tools, known_uris, known_templates, \
                template_regexes, known_prompts = \
                process_server_capabilities(
                    tools_resp, resources_resp,
                    templates_resp, prompts_resp
                )

            print_capabilities(known_uris, known_templates, known_prompts)
            await chat_loop(
                session, ollama_tools, known_uris,
                template_regexes, known_prompts
            )
\end{lstlisting}

\subsection{Capability discovery}
On startup the client issues four parallel MCP calls (\texttt{list\_tools}, \texttt{list\_resources}, \texttt{list\_resource\_templates}, \texttt{list\_prompts}) and hands the results to \texttt{process\_server\_capabilities} (defined in \texttt{client\_helper.py}). The outputs are:

\begin{itemize}
  \item \texttt{ollama\_tools} — tool definitions translated into the JSON schema format Ollama expects.
  \item \texttt{known\_uris} — set of static resource URIs for \texttt{@mention} resolution.
  \item \texttt{template\_regexes} — compiled regexes for URI template matching.
  \item \texttt{known\_prompts} — set of prompt names for \texttt{/command} detection.
\end{itemize}

\subsection{The chat loop}

\begin{lstlisting}[caption={chat\_loop}]
async def chat_loop(session, ollama_tools, known_uris,
                    template_regexes, known_prompts):
    messages = []
    while True:
        try:
            user_input = input(c("You: ", BOLD, BLUE)).strip()
        except (EOFError, KeyboardInterrupt):
            print(c("\nGoodbye!", DIM))
            return

        if user_input.lower() in ("exit", "quit"):
            print(c("Goodbye!", DIM))
            return

        if not user_input:
            continue

        turn = await build_user_turn(
            session, user_input, known_uris,
            template_regexes, known_prompts
        )
        messages.extend(turn)
        await run_ollama_turn(session, messages, ollama_tools)
\end{lstlisting}

The loop maintains the full \texttt{messages} list across turns, giving the model conversation memory. Each iteration:
\begin{enumerate}
  \item Reads a line of user input.
  \item Calls \texttt{build\_user\_turn} to resolve any special syntax.
  \item Calls \texttt{run\_ollama\_turn} to get a model response.
\end{enumerate}

\subsection{Building a user turn}

\begin{lstlisting}[caption={build\_user\_turn}]
async def build_user_turn(session, user_input, known_uris,
                           template_regexes, known_prompts):
    # /prompt_name key=val  ->  fetch prompt messages from server
    parsed = parse_prompt_command(user_input, known_prompts)
    if parsed is not None:
        prompt_name, arguments = parsed
        result = await session.get_prompt(
            prompt_name,
            arguments=arguments if arguments else None
        )
        return [
            {"role": msg.role, "content": msg.content.text}
            for msg in result.messages
            if hasattr(msg.content, "text")
        ]

    # @uri mentions  ->  fetch resource content and prepend it
    parts = []
    for uri, is_template in parse_resource_mentions(
        user_input, known_uris, template_regexes
    ):
        result = await session.read_resource(uri)
        content_text = "\n".join(
            block.text for block in result.contents
            if hasattr(block, "text")
        )
        parts.append(f"Resource {uri}:\n{content_text}")

    resource_context = "\n\n".join(parts) if parts else None
    content = (
        f"{user_input}\n\n{resource_context}"
        if resource_context else user_input
    )
    return [{"role": "user", "content": content}]
\end{lstlisting}

Two special syntaxes are handled before sending input to the model:

\begin{description}[leftmargin=3em]
  \item[\texttt{/prompt\_name [key=val \ldots]}] Fetches the named prompt from the server and returns its messages directly, bypassing the normal user message. The messages are injected into the conversation as-is (with their original \texttt{role} values).
  \item[\texttt{@uri}] Fetches the resource at \texttt{uri} (or a URI that matches a template) and prepends its text content to the user message.
\end{description}

\subsection{Running an Ollama turn}

\begin{lstlisting}[caption={run\_ollama\_turn}]
async def run_ollama_turn(session, messages, ollama_tools):
    while True:
        response = ollama.chat(
            model=MODEL,
            messages=messages,
            tools=ollama_tools or None
        )
        assistant_message = response.message
        messages.append(assistant_message)

        if not assistant_message.tool_calls:
            print(c("Assistant: ", BOLD, GREEN) + assistant_message.content)
            return

        for tool_call in assistant_message.tool_calls:
            tool_name = tool_call.function.name
            tool_args = dict(tool_call.function.arguments)
            result = await session.call_tool(tool_name, tool_args)
            result_text = "\n".join(
                block.text for block in result.content
                if hasattr(block, "text")
            )
            messages.append({"role": "tool", "content": result_text})
\end{lstlisting}

This function implements the standard \emph{agentic loop}:
\begin{enumerate}
  \item Send the current message history to the model.
  \item If the model replies with plain text and no tool calls, print and return.
  \item If the model requests tool calls, execute each one via the MCP session, append the result as a \texttt{tool} message, and loop back to step 1.
\end{enumerate}

% ══════════════════════════════════════════════════════════════════════════════
\section{Adding New Capabilities}
% ══════════════════════════════════════════════════════════════════════════════

\subsection{Adding a tool}

Create a new \texttt{.py} file anywhere inside \texttt{tools/}. Import the shared \lstinline{mcp} instance and decorate your function with \lstinline{@mcp.tool()}. The server will discover and register it automatically on the next start.

\begin{lstlisting}[caption={tools/my\_new\_tool.py}]
from mcp_instance import mcp


@mcp.tool()
async def my_new_tool(param_one: str, param_two: int = 0) -> dict:
    """One-line description shown to the model.

    Longer explanation if needed. The model uses both the function
    signature and this docstring to decide when and how to call the tool.

    Args:
        param_one: Description of the first parameter.
        param_two: Description of the second parameter.

    Returns:
        A dictionary with the result.
    """
    return {"result": f"processed {param_one} with offset {param_two}"}
\end{lstlisting}

\paragraph{Guidelines.}
\begin{itemize}
  \item Use \texttt{async def}; the server is fully async.
  \item Type-annotate all parameters — FastMCP generates the JSON schema from annotations.
  \item Write a descriptive docstring. The model reads it to decide when to call the tool.
  \item Return a plain Python value (dict, str, list, etc.); FastMCP serialises it.
\end{itemize}

\subsection{Adding a resource}

Create a new \texttt{.py} file inside \texttt{resources/}. Use a URI string as the decorator argument. The URI can contain \texttt{\{param\}} placeholders for templated resources.

\begin{lstlisting}[caption={resources/my\_static\_resource.py — fixed URI}]
from mcp_instance import mcp


@mcp.resource("myapp://config/defaults")
async def get_defaults() -> str:
    """Return default application configuration."""
    return "timeout=30\nretries=3\nlog_level=INFO"
\end{lstlisting}

\begin{lstlisting}[caption={resources/my\_template\_resource.py --- URI template}]
from mcp_instance import mcp


@mcp.resource("myapp://items/{item_id}")
async def get_item(item_id: str) -> str:
    """Fetch a specific item by ID.

    Args:
        item_id: The item identifier extracted from the URI.
    """
    # item_id is automatically extracted from the URI path
    return f"Data for item: {item_id}"
\end{lstlisting}

\paragraph{Accessing a resource from the client.}
Users type \texttt{@myapp://config/defaults} or \texttt{@myapp://items/42} in the prompt. The client resolves it against the known URIs and templates, fetches the content, and injects it into the message.

\subsection{Adding a prompt}

Create a new \texttt{.py} file inside \texttt{prompts/}. A prompt function returns a string (or a list of message dicts for multi-turn templates). Its parameters become the keyword arguments users can pass with \texttt{/prompt\_name key=val}.

\begin{lstlisting}[caption={prompts/my\_prompt.py}]
from mcp_instance import mcp


@mcp.prompt()
async def summarise_document(text: str, style: str = "bullet") -> str:
    """Summarise a block of text.

    Args:
        text:  The document to summarise.
        style: Output style - "bullet" or "paragraph".
    """
    return (
        f"Summarise the following text in {style} form.\n\n"
        f"Text:\n{text}"
    )
\end{lstlisting}

\paragraph{Invoking from the client.}
\begin{lstlisting}[language=bash, style=python, numbers=none]
/summarise_document text="Hello world." style=bullet
\end{lstlisting}

The client detects the \texttt{/} prefix, matches the name against known prompts, and calls \texttt{session.get\_prompt} with the parsed keyword arguments.

% ══════════════════════════════════════════════════════════════════════════════
\section{Existing Modules Reference}
% ══════════════════════════════════════════════════════════════════════════════

\subsection{Tools}

\subsubsection{\texttt{verify\_product\_transparency} (tools/coke\_product\_transparency.py)}

Looks up a product code in an in-memory database and returns detailed transparency information: ingredients, sourcing, nutritional facts, sustainability metrics, and certifications. Returns a not-found message if the code is unknown, which can signal a counterfeit.

\textbf{Parameter:} \texttt{product\_code: str} — e.g.\ \texttt{"COKE-355ML-001"}.

\subsubsection{\texttt{query\_crisis\_protocol} (tools/coke\_crisis\_protocol.py)}

Looks up a named crisis management protocol from an in-memory dictionary. Returns priority, step-by-step action list, and contact addresses. Returns the list of available protocol names when no match is found. The comments in the file illustrate how a real implementation would replace the dictionary with a RAG vector search.

\textbf{Parameter:} \texttt{crisis\_type: str} — e.g.\ \texttt{"product recall"}, \texttt{"contamination"}.

\subsection{Resources}

\subsubsection{\texttt{coke://media/advertisements/\{ad\_id\}} (resources/coke\_advertisements.py)}

URI template resource. Returns a URL and short description for a named advertisement asset. Example IDs: \texttt{hilltop-1971}, \texttt{polar-bears-1993}.

\subsubsection{\texttt{coke://locations/mit-campus} (resources/coke\_campus\_locations.py)}

Static resource. Reads and returns \texttt{resources/data/mit\_campus.txt} — a directory of Coca-Cola vending and dining locations on the MIT campus.

\subsubsection{\texttt{coke://regulations/massachusetts/\{regulation\_id\}} (resources/coke\_regulations.py)}

URI template resource. Reads \texttt{resources/data/massachusetts\_regulations.json} and returns the entry matching \texttt{regulation\_id} as formatted JSON. Example IDs: \texttt{bottle-bill}, \texttt{sugar-tax}, \texttt{labeling}, \texttt{school-sales}.

\subsection{Prompts}

\subsubsection{\texttt{evaluate\_beverage\_choice} (prompts/coke\_evaluate\_beverage.py)}

Returns a detailed system prompt that guides the model through a five-step chain-of-thought for beverage selection: understanding the user's needs, evaluating options transparently, considering trade-offs, giving an honest recommendation, and educating on long-term patterns.

\textbf{Parameter:} \texttt{beverage\_context: str} — e.g.\ \texttt{"choosing a drink for lunch"}.

% ══════════════════════════════════════════════════════════════════════════════
\appendix
% ══════════════════════════════════════════════════════════════════════════════

\section{Client Helper (\texttt{client\_helper.py})}
\label{appendix:client_helper}

\texttt{client\_helper.py} contains all parsing logic and terminal formatting used by the client. It has no side effects at import time and carries no dependency on the MCP server.

\subsection{ANSI colour helpers}

\begin{lstlisting}[caption={Colour constants and the c() helper}]
RESET  = "\033[0m"
BOLD   = "\033[1m"
DIM    = "\033[2m"
CYAN   = "\033[36m"
GREEN  = "\033[32m"
YELLOW = "\033[33m"
RED    = "\033[31m"
BLUE   = "\033[34m"

def c(text, *codes):
    return "".join(codes) + str(text) + RESET
\end{lstlisting}

\texttt{c(text, *codes)} concatenates any number of ANSI escape codes before \texttt{text} and appends the reset code after it. Usage: \lstinline{c("hello", BOLD, GREEN)}.

\subsection{URI template conversion}

\begin{lstlisting}[caption={uri\_template\_to\_regex}]
def uri_template_to_regex(template):
    """Convert a URI template like 'myapp://data/{id}'
    to a compiled regex."""
    escaped = re.escape(template)
    pattern = re.sub(
        r'\\\{(\w+)\\\}',
        r'(?P<\1>[^/]+)',
        escaped
    )
    return re.compile(f'^{pattern}$')
\end{lstlisting}

Converts an RFC 6570-style URI template (e.g.\ \texttt{coke://regulations/massachusetts/\{regulation\_id\}}) into a compiled Python regex with named capture groups. This allows the client to check whether a user-typed \texttt{@uri} matches any template, without needing the MCP server to enumerate all possible parameterised values.

\subsection{Capability processing}

\begin{lstlisting}[caption={process\_server\_capabilities and build\_tools\_for\_ollama}]
def process_server_capabilities(tools_resp, resources_resp,
                                 templates_resp, prompts_resp):
    ollama_tools = build_tools_for_ollama(tools_resp.tools)
    known_uris = {str(r.uri) for r in resources_resp.resources}
    known_templates = [
        str(t.uriTemplate) for t in templates_resp.resourceTemplates
    ]
    template_regexes = [
        uri_template_to_regex(t) for t in known_templates
    ]
    known_prompts = {p.name for p in prompts_resp.prompts}
    return (ollama_tools, known_uris, known_templates,
            template_regexes, known_prompts)


def build_tools_for_ollama(tools):
    return [
        {
            "type": "function",
            "function": {
                "name": tool.name,
                "description": tool.description or "",
                "parameters": tool.inputSchema,
            },
        }
        for tool in tools
    ]
\end{lstlisting}

\texttt{process\_server\_capabilities} transforms the four raw MCP list responses into the data structures the client loop needs. \texttt{build\_tools\_for\_ollama} re-packages MCP tool descriptors into Ollama's \texttt{function} schema, which the Ollama Python library passes directly to the model.

\subsection{Input parsing}

\begin{lstlisting}[caption={Regex patterns and parsers}]
RESOURCE_PATTERN = re.compile(r'@(\S+)')
PROMPT_PATTERN   = re.compile(r'/(\w+)((?:\s+\w+=\S+)*)\s*$')
KWARG_PATTERN    = re.compile(r'(\w+)=(\S+)')


def parse_prompt_command(user_input, known_prompts):
    match = PROMPT_PATTERN.match(user_input.strip())
    if not match or match.group(1) not in known_prompts:
        return None
    prompt_name = match.group(1)
    arguments = {k: v for k, v in KWARG_PATTERN.findall(match.group(2))}
    return prompt_name, arguments


def parse_resource_mentions(user_input, known_uris, template_regexes):
    valid = []
    for uri in RESOURCE_PATTERN.findall(user_input):
        if uri in known_uris:
            valid.append((uri, False))
        elif any(r.match(uri) for r in template_regexes):
            valid.append((uri, True))
        else:
            print(c(f"  x unknown resource: {uri}", RED))
    return valid
\end{lstlisting}

\paragraph{\texttt{parse\_prompt\_command}.}
Matches input of the form \texttt{/name [key=val \ldots]} against the known prompt names. Returns \texttt{(name, \{key: val, \ldots\})} on a match, or \texttt{None} if the input is not a prompt command or the name is not registered.

\paragraph{\texttt{parse\_resource\_mentions}.}
Finds all \texttt{@token} substrings in the input using \texttt{RESOURCE\_PATTERN}. Each token is checked first against the set of exact static URIs, then against every compiled template regex. Unrecognised URIs are printed as errors in red. Returns a list of \texttt{(uri, is\_template)} pairs.

\subsection{\texttt{print\_capabilities}}

\begin{lstlisting}[caption={print\_capabilities}]
def print_capabilities(known_uris, known_templates, known_prompts):
    print(c("MCP Chat", BOLD, CYAN)
          + c("  -  type 'exit' or 'quit' to stop", DIM))
    print()
    print(c("Usage:", BOLD))
    print(c("  @<uri>", CYAN)
          + "                        attach a resource")
    print(c("  @<uri-template-with-values>", CYAN)
          + "   attach a resource template")
    print(c("  /<prompt>", CYAN)
          + c(" [key=val ...]", DIM)
          + "        inject a prompt")
    print()
    if known_uris:
        print(c("  Resources:          ", BOLD)
              + c(', '.join(sorted(known_uris)), CYAN))
    if known_templates:
        print(c("  Resource templates: ", BOLD)
              + c(', '.join(sorted(known_templates)), CYAN))
    if known_prompts:
        print(c("  Prompts:            ", BOLD)
              + c(', '.join(sorted(known_prompts)), CYAN))
    print()
\end{lstlisting}

Prints a usage banner with a live listing of all available resources, templates, and prompts discovered from the server. Shown once at startup so the user knows what they can reference.

\end{document}
